\section{Estimate Tab}
\label{sec:estimate}

The \app{Estimate} tab predicts beam time before a measurement. Given
an element, an electronic transition, an isotope chain, and beam
parameters, it computes per-peak signal rates, the on-peak time
required to reach a target statistical significance, the number of
voltage sweeps needed to cover the requested scan range, and a
grand-total run time. A \app{Spectrum only} switch suppresses the
timing pass and produces only the simulated lineshape --- useful when
planning the scan window itself.

The tab is a single surface with three resizable columns: the global
parameters on the left, the isotope list in the middle with the
\app{Run Options} group under it, and a right column stacking the plot
display, the \app{Voltage Range Estimator}, and the streaming log. The
expand button on the plot panel temporarily expands the plots over the
whole tab; click it again to bring the other panels back.

\subsection{Global parameters}

The left column holds settings shared by every isotope.

\begin{description}
    \item[\app{Element \& Transition}.]
        Element symbol, atomic number $Z$, the lower and upper levels
        of the electronic transition in cm$^{-1}$, and their total
        angular momenta $J_l$, $J_u$.
    \item[\app{Beam}.]
        Acceleration voltage (kV), ion charge state, and laser-ion
        geometry --- \app{collinear} or \app{anti-collinear}.
    \item[\app{Laser}.]
        Harmonic order and the setpoint of the \emph{fundamental}
        laser in cm$^{-1}$. As an alternative to a fixed setpoint,
        tick \app{Anchor to isotope}, choose an isotope and its state
        (g.s.\ or isomer), and enter the voltage offset \app{dV} from
        $V_0$ at which that resonance should sit; \denis{} then
        back-solves the laser frequency required to place the
        centroid at that voltage.
    \item[\app{Linewidth}.]
        Full-width-half-maximum used when drawing simulated spectra
        (MHz). Does not affect timing.
    \item[\app{Scan}.]
        Voltage range (full window, $\pm$ half on either side of
        $V_0$), step size, and dwell time per step (ms).
\end{description}

The checkable \app{Statistics (Timing Mode)} group is the master
switch for timing: unchecking it puts the whole tab in spectrum-only
mode (rates, backgrounds, and yields are hidden, and the scan step
size and dwell time are disabled). When checked it exposes:

\begin{itemize}
    \item \app{Required $\sigma$} --- target statistical significance per
        peak; defaults to $3\sigma$.
    \item \app{Uniform timing (global default)} --- when ticked, every
        measured peak gets the same on-peak time, computed for the
        strongest peak assuming full yield. Useful when a single dwell
        setting must cover the whole scan. Any isotope can override
        this locally with its own \app{Uniform timing} control.
\end{itemize}

\subsection{Isotope entries}

The middle column contains one expandable panel per isotope, added
with the \app{+ Add Isotope} button at the top. Panels are reordered
by dragging the hamburger handle and are color-striped by position so
the same isotope reads consistently across the tab. Each panel
exposes:

\begin{description}
    \item[Mass number $A$ and label] used in plot legends and log
        headers.
    \item[Nuclear spin $I$] integer or half-integer.
    \item[$\mu$ ($\mu_N$) and $Q$ (b)] magnetic dipole and electric
        quadrupole moments.
    \item[Isotope shift (MHz)] centroid offset from the reference
        isotope.
    \item[Yield (ions/s)] expected ion rate (a scientific-notation
        field).
    \item[Efficiency] one detected photon per
        $1/\eta$ ions; the log echoes this expanded as e.g.\
        ``$1$ photon per $500$ ions''.
    \item[Peaks to measure] how many of the strongest hyperfine peaks
        enter the timing estimate (default 1; raise it when several
        components are diagnostic).
    \item[Background] either a flat rate (Hz), or a continuous rate
        combined with a TOF-gate width (\si{\micro\second}) and the
        dwell time, which then collapses to an effective rate
        automatically.
\end{description}

Exactly one isotope is marked as the \app{Reference}. Its hyperfine
$A$ and $B$ constants (entered in MHz, lower- and upper-state
separately) are used to scale the constants of every other isotope
via the Schmidt-moment ratio. Tick \app{Override A/B constants} on a
non-reference isotope to bypass scaling and supply explicit values ---
useful when high-quality measured constants are already available.

\paragraph{Schmidt-moment auto-calculation.}
For an isotope whose $\mu$ is unknown, tick \app{Use Schmidt moment}
to derive it from the shell model. The single-particle case takes a
nucleon type (\code{p}/\code{n}) and an orbital (e.g.\ \code{f7/2}).
For odd-odd nuclei, choose the two-particle variant and supply two
orbital configurations together with the resulting spin $I$. The
unquenched Schmidt $\mu$ is auto-filled into the panel's $\mu$ field
(undoable with \code{Ctrl+Z}); the configured $g_s$ quenching factor
is applied during the run, and the log prints the configuration
string next to the value actually used. The same calculator is
available standalone under \menu{Tools $\rightarrow$ Schmidt Moment
Calculator} (Section~\ref{sec:tools}).

\paragraph{Isomers.}
Tick \app{Has isomer} to expand a sub-panel for an isomeric state.
The isomer carries its own spin, moments, isotope shift, and one of
three yield specifications: an explicit \app{Yield} (independent of
the ground state), a \app{Production ratio} in $[0,1]$ that splits
the parent's yield, or a \app{Spin distribution} $\sigma$ that
estimates the split from spin statistics --- the latter is useful for
surface-ion sources where high-spin states are populated
preferentially. The isomer contributes its own peaks to the simulated
spectrum and its own row to the timing table.

\subsection{Running the estimation}

Click \app{Run} in the \app{Run Options} group (under the isotope
list). The plots, the log, and the voltage range estimator stay empty
until then. From the first run on, every change (an edit, an undo, an
estimator commit, the palette, \app{Generate plots}) re-runs the
estimation about half a second after the last edit and updates all
three; an incomplete configuration is reported in the status bar and
the log, and the last good plots stay up. The log contains, in order:

\begin{enumerate}
    \item A header summarising the inputs and, in anchor mode, the
        derived laser setpoint.
    \item Per isotope: signal rate, on-peak time per measured peak,
        the number of sweeps required to accumulate enough statistics
        on each peak, and a subtotal in hours and in 8\,h shifts.
    \item A \app{Nuclear parameters} summary table with the $\mu$
        values actually used (config or Schmidt) and the resulting
        $A$, $B$ constants.
    \item The grand-total run time across all isotopes.
    \item Two peak lists --- one sorted by intensity, one by voltage
        --- covering every hyperfine component of every isotope and
        isomer.
\end{enumerate}

Runs go to a temporary folder that is deleted when DENIS closes.
\app{Save} copies the latest run into the chosen \app{Output directory}
as a \code{cls\_<timestamp>} folder: the log written verbatim to a
\code{.log} file, the plot PDFs, the plot arrays and peak table
(\code{estimate\_results.npz}, \code{peaks.csv}), and every input the
run used (\code{estimate\_config.yaml}: the parameter boxes, the
isotopes and the reference). The session save
file records the saved folder --- so a saved session restores its
plots, peak table, and log on load without re-running.

\subsection{Simulated spectra}

The plot panel shows a simulated voltage spectrum per isotope on the
configured scan range, with vertical dashed lines at each hyperfine
peak position. The
\app{View} drop-down switches between \app{Individual Spectra} (one
panel per isotope), \app{Combined Overview} (all isotopes overlaid),
and \app{Peak List} (a sortable table of every hyperfine component).
\app{Edit Plot} opens the matplotlib style editor; \app{Open PDF}
reveals the exported PDF in the system file viewer.

\subsection{Voltage range estimator}

The \app{Voltage Range Estimator} picks a laser setpoint and a sweep
interactively, using the same physics as \app{Run}. Its map
plots, against the fundamental setpoint, the voltage offset $\Delta V$
at which every hyperfine component resonates: one line per component,
solid for ground states and dashed for isomers, in each isotope's
colour. A red bar marks the current setpoint, with a red ring where it
crosses a selected peak's line inside the sweep, fading across the
margin towards the sweep's edges, where scanning is riskier (the readout
calls those crossings \emph{edge} and counts them by weight in its
coverage line); a shaded band marks the $\pm$\app{dV} sweep; thin
\emph{centre} markers show the setpoint that puts each state's centroid
at $\Delta V = 0$. Beside the map, the spectrum at the current setpoint
shares the $\Delta V$ axis, and a readout lists the setpoint, laser
frequency, beam voltage and every component's offset, $\Delta V$,
intensity, and whether it lies inside the sweep.

Drag the red bar, or click the map and use the arrow keys
($\pm$0.001\,cm$^{-1}$, with Shift $\pm$0.01\,cm$^{-1}$); clicking a
centre marker snaps to it. The mouse wheel zooms the map around the
cursor, dragging away from the bar pans it, and a double-click refits
it; the zoom is kept while you adjust the setpoint and the sweep. When you finish (mouse release, Enter in the
\app{Setpoint} box, or a pause after the arrow keys), the setpoint is
written to \app{Laser} $\rightarrow$ \app{Setpoint}; in anchor mode the
\app{Setpoint} radio is selected and the anchor fields keep their
values. Entering \app{dV} writes \app{Scan} $\rightarrow$
\app{Voltage range} $= 2\,$\app{dV}. Both writes can be undone.
\app{Peaks} picks the components the sweep must cover: by default
each state's strongest \app{Peaks to measure} (the ones the run time
counts), all components, or any ticked per state; selected components
are drawn bold and marked in the readout, whose \emph{Peaks} line
gives the DAQ range [start, stop] that covers them at the current
setpoint, with 100\,V of baseline on each side (the gear button beside
\app{Optimise} changes this), start and stop rounded outward to whole
steps. \app{Optimise} sets \app{Scan} $\rightarrow$ \app{Voltage range}
to that range's width (one undo step); the setpoint is left for you to
place. With \app{Show} on a single state it covers that state alone.
The same gear holds a \app{Readout offset}: a signed shift in volts for
a known DAQ readout shift (negative moves the peaks to lower
$\Delta V$). The component lines, rings, spectrum and sweep band move by
it, and the readout's $\Delta V$ values and scan range include it, so
they can be entered as DAQ scanning ranges directly. The axes do not
move, and the hyperfine model, the plots and the log are unchanged.
\app{Window} only sets the plotted $\Delta V$ range (\app{auto} fits
the components), and \app{Show} limits the view to one state. The
estimator appears after the first \app{Run}; its commits then re-run
the estimation like any other change.

\subsection{Loading a previous estimation}

\app{Load\dots} (next to \app{Run} and \app{Save}) reopens any
saved \code{cls\_<timestamp>} run folder for viewing --- plots,
peak table, and log are restored without re-running and without
creating new output --- and its inputs are restored from
\code{estimate\_config.yaml}, so \app{Run} repeats it straight away.
A loaded run, like a loaded configuration, stays a static view until
\app{Run} is pressed; folders saved before the inputs were kept load
their outputs only. Runs from before the \code{.npz} persistence
existed hold only PDFs; their pages are rendered onto the live
matplotlib canvas, so pan/zoom/save still work.
